
%%%%%%%%%%%%%%%%%%%%%%%%%%%%%%%%%%%%%%%%%%%%%%%%%%%%%%%%%%%%%%%%%%%%%%%%%%%%%%%%
\section{General info about the course and Introduction}
\sectionframe{1}{General info about the course and Introduction}

%--------------------------
\myframe{Schedule}{
    \begin{tabular}{c|c|l}
        \textbf{Date} & \textbf{Time} & \textbf{Tentative Topic}\\
        Mon. 27/04/2026 & h10.30-12:30 (2h) & Intro e Preliminaries\\
        Mon. 04/05/2026 & h09.30-12:30 (3h) & Preliminaries/Linear Theory\\
        Mon. 11/05/2026 & h09.30-12:30 (3h) & Linear Theory\\
        Mon. 18/05/2026 & h09.30-13:30 (4h) & Linear Theory + Essay\\
        Mon. 25/05/2026 & h09.30-13:30 (4h) & Linear/Nonlinear Theory + Essay\\
        Mon. 01/06/2026 & h09.30-13:30 (4h) & Nonlinear Theory/Adds-on + Essay 
    \end{tabular}

    \vspace{.5cm}
    \pause
    \begin{tabular}{c|c|l}
        \textbf{Date} & \textbf{Time} & \textbf{Tentative Topic}\\
        Mon. 27/04/2026 & h10.30-12:30 (2h) & Intro e Preliminaries\\
        \sout{\textcolor{red}{Mon. 04/05/2026}} & \sout{\textcolor{red}{h09.30-12:30 (3h)}} & Preliminaries/Linear Theory\\
        Fri. 08/05/2026 & h10.30-13:30 (3h) & \\
        Mon. 11/05/2026 & h09.30-12:30 (3h) & Linear Theory\\
        \sout{\textcolor{red}{Mon. 18/05/2026}} & h09.30-13:30 (4h) & Linear Theory + Essay\\
        Fri. 15/05/2026 &  & \\
        \sout{\textcolor{red}{Mon. 25/05/2026}} & h09.30-13:30 (4h) & Linear/Nonlinear Theory + Essay\\
        Fri. 22/05/2026 & h09.30-13:30 (4h) & Linear/Nonlinear Theory + Essay\\
        \sout{\textcolor{red}{Mon. 01/06/2026}} & h09.30-13:30 (4h) & Nonlinear Theory/Adds-on + Essay\\
        Mon. 25/05/2026 & h09.30-13:30 (4h) & Nonlinear Theory/Adds-on + Essay 
    \end{tabular}
}

%--------------------------
\myframe{Exam}{
    % \begin{enumerate}
    %     \setcounter{enumi}{\numexpr1-1\relax}
    %     \item Proofs of theorems seen during lessons 
    %     \begin{itemize}
    %         \item during lessons we will encounter some (few) basic
    %         theorems/corollaries/properties. You are asked \underline{\textbf{to
    %         try}} to prove them formally on your own.
    %         \item do your best and try not to cheat!
    %         \item you can work in a group if preferred. 
    %     \end{itemize}
    % \end{enumerate}
    % \vspace{0.2mm}
    % AND
    % \vspace{0.2mm}
    % \begin{enumerate}
    %     \setcounter{enumi}{\numexpr2-1\relax}
    %     \item Practical coding essay:
    %     \begin{itemize}
    %         \item on a selected application: pick a topic that is relatable to
    %         your need/research.
    %         \item partially developed during class.
    %         \item programming language: python (MATLAB is an option but try to
    %         avoid it), either using jupyter notebook or script based. Make sure all
    %         necessary dependencies to run are available.
    %         \item no use of predefined libraries with KF implementations.
    %         \item you can work in a group if preferred.
    %         \item \underline{\textbf{objective}:} explore (possibly all) the
    %         topics covered during the course, from linear to non-linear KF. You
    %         might need to start from modeling assumptions, linearize if
    %         necessary, discuss noise statistics and then implement the filter.
    %         Try to explore limits and nuances as if it were a small research
    %         project.
    %     \end{itemize}
    % \end{enumerate}

    Practical coding essay + report:
        \begin{itemize}
            \item on a selected application: pick a topic possibly related to
            your need/research.
            \item partially developed during class.
            \item programming language: python -- either using jupyter notebooks
            (recommended) or script based. Make sure all necessary dependencies
            to run are available.\\\textbf{For other programming language ask first!!}
            \item no use of predefined libraries with KF implementations.
            \item you can work in a group if preferred.
            \item \underline{\textbf{objective}:} explore (possibly all) the
            topics covered during the course, from linear to non-linear KF. You
            might need to start from modeling assumptions, linearize if
            necessary, discuss noise statistics and then implement the filter.
            Try to explore limits and nuances as if it were a small research
            project.
            \item present results in a pdf report\footnote{For well organized
            runnable Jupyter notebooks this is optional. If, upon agreement,
            another language is used this is mandatory.}
        \end{itemize}

}

%--------------------------
\myframe{References and additional material}{
    \begin{itemize}
        \item [IT] \textbf{Picci, G. (2007). Filtraggio statistico (Wiener, Levinson, Kalman) e applicazioni. Libreria Progetto.}
        \item [EN] Sanso F., Pileggi, C. and Biagi, L. (2024). A tutorial on Kalman filtering. The linear, extended and unscented filters. (pdf available as resource in the material folder)
        \item [EN] Anderson, B. D., and Moore, J. B. (2005). Optimal filtering. Courier Corporation.
    \end{itemize}

    \vspace{5mm}

    G-drive:\\
    {\small
    \url{https://drive.google.com/drive/folders/1BbuSTiO4wym0qOAqeLUkiVum_14kRRzD}}

    \vspace{5mm}

    github: for complete slides, code and additional resources\\
    {\small
    \url{https://github.com/MarcoTodescato/linear_and_nonlinear_kalman_filtering}}

    \vspace{5mm}
    For Q/A and discussion we can set up a Discord channel
}


%--------------------------
\myframe{Goal}{%
How to formulate estimation problems encountered in many engineering fields and
applications as instances of a general dynamical system state estimation problem
and how to derive the mathematical solution of the latter problem (possibly
depending on additional assumption on the functional family of the considered
dynamical system).%

\vspace{10mm}

\pause
\textbf{Essentially, this will lead to the well-known (linear) Kalman filter and
its (nonlinear) extensions.}%
}

%--------------------------
\myframe{A bit of history: the theories of statistical filtering}{%

\begin{exampleblock}{Wiener-Kolmogorov (1940)}
    \begin{itemize}
        \item $I=(-\infty, t]$, infinite time support
        \item $\{\x(t)\}$, $\{\y(t)\}$ are weakly jointly stationary
    \end{itemize}
    The impulse response of the optimal filter is time-invariant hence it can be
    analytically computed once and for all times, knowing the processes
    statistics.
\end{exampleblock}

\pause

\begin{block}{Levinson (1950)}
    \begin{itemize}
        \item no more infinite (lower unbounded) time support
    \end{itemize}
    The impulse response of the filter is time-varying but can be efficiently
    computed at time $t$ using a recursive structure which exploits the
    computations until $t-1$. 
\end{block}

\pause

\begin{alertblock}{Kalman (1960)}
    \begin{itemize}
        \item signals can be modeled as stochastic dynamical systems of finite
        size
    \end{itemize}
    Looks for recursive finite-memory schemes that resembles that same dynamical
    system structure of finite size.
\end{alertblock}
}

%--------------------------
\myframe{A bit of history: a change of perspective}{%
Revolutionary apporach introduced by R.E. Kalman in 1960 to the problem of
statistical filtering.%

\vspace{5mm}

Change of perspective:
\begin{itemize}
    \item \textbf{Wiener-Kolmogorov}: explicit computation of the filter
    \textbf{transfer function}, practically useful only for signals
    characterized by a \textbf{rational} spectral density
        \begin{itemize}
            \item implicit computations possible only for rational signals
            \item actual physical implementation of the filter as a set of
            difference (recursive) equations only for rational representations
        \end{itemize}

    \pause
    \medskip
    \item \textbf{Kalman}: hypothesis that the signals can be described as
    (linear) \textbf{dynamical systems} of finite dimension
        \begin{itemize}
            \item no transfer functions involved
            \item natural formulation of the filter as a fixed algorithmic
            structure essentially consisting of a discrete-time (linear)
            dynamical system of finite memory, driven by measurements
            \item implementation always possible as a recursive algorithm 
        \end{itemize}
\end{itemize}
}

%--------------------------
\myframe{Prediction vs Filtering vs Smoothing}{%
Assume $\{\x(t)\}$ and $\{\y(t)\}$ are (second order\footnote{First and second
moment statistics exist and are finite.}) processes of known dimensions. Denote
with $\xhat(t|I)$ the estimator of $\x(t)$ given $\{\y(t)\}_{t\in I}$, $I$ being
the time support.

\vspace{.2cm}

It is possible to identify 3 classes of \underline{\textbf{estimation problems}}.

\vspace{.2cm}

\pause
\begin{block}{Prediction}
    Find $\xhat(t+h|I)$ being $I=[t_0,t]$ and $h>0$, that is the estimator
    $h$-step ahead in the future from measurements up to time $t$.
\end{block}

\pause
\begin{block}{Filtering}
    Find $\xhat(t|I)$ being $I=[t_0,t]$, that is the estimator at current time
    $t$ from all available measurements up to $t$. 
\end{block}

\pause
\begin{block}{Smoothing}
    Find $\xhat(t|I)$ being $I=[t_0,t_1]$ with $t<t_1$, that is the estimator at
    time $t$ from all available measurements up to $t_1$ in the future.
\end{block}
}
